% Appendices A-C of the thesis. Needs: booktabs, amsmath.
% Tables in C are generated by scripts/thesis_tables.py (thesis/tables/*.tex).

\appendix

% ---------------------------------------------------------------------------
\chapter{Experimental Setup and Hyperparameters}
\label{app:setup}

This appendix lists the settings behind every run in Chapter~6, so that each
experiment can be reproduced. All arms share the baseline configuration of
Section~\ref{app:setup_vae}. Each arm changes a single setting relative to that
baseline (Table~\ref{tab:app_arms}).

\section{Base model}
\label{app:setup_base}

Both stages start from LTX-Video 2B (v0.9.6, \texttt{dev}). Its causal video
VAE compresses by $32\times$ in each spatial dimension and $8\times$ in time,
into 128 latent channels. A 101-frame simulation is padded to $105 = 8k+1$
frames. At $256\times256$ this gives a latent of $14\times8\times8$ with 128
channels. The pretrained VAE expects three (RGB) input and output channels. The
first convolution of the encoder (\texttt{conv\_in}) and the last convolution
of the decoder (\texttt{conv\_out}) are therefore widened from three to four
channels. The channel order is density, $x$-momentum, $y$-momentum and
pressure, so pressure occupies the new fourth slot. The weights for this slot
are initialised to zero, and the three pretrained slots keep their weights.

\section{VAE fine-tuning}
\label{app:setup_vae}

\begin{table}[htbp]
\centering
\small
\caption{Shared VAE fine-tuning configuration (the baseline arm).}
\label{tab:app_vae_hparams}
\begin{tabular}{ll}
\toprule
Setting & Value \\
\midrule
Trainable modules & entire encoder (incl.\ widened \texttt{conv\_in}) and decoder \\
Loss & $\mathcal{L} = 1.0\,\mathcal{L}_\mathrm{RMSE} + 50.0\,\mathcal{L}_{H^1}$ \\
Optimiser & AdamW, weight decay 0 \\
Learning rate & $5\times10^{-5}$, 200 linear warm-up steps (from $0.01\times$), cosine decay to $10^{-6}$ \\
Epochs & 20 (4000 training simulations per epoch) \\
Batch size & 32 (4 per device) \\
Gradient clipping & global norm 5.0 \\
Precision & bf16 mixed precision, gradient checkpointing \\
Checkpoint used & best validation VRMSE (epoch 20 for every arm) \\
Seed & 42 \\
\bottomrule
\end{tabular}
\end{table}

$\mathcal{L}_\mathrm{RMSE}$ is the root-mean-square error of the normalised
fields. $\mathcal{L}_{H^1}$ is the squared difference of their first-order
spatial gradients, $\lVert \nabla \hat u - \nabla u \rVert^2$. The
regularisation arms add one of two further terms:

\paragraph{KL divergence.} The standard VAE prior term on the encoder
posterior $q(z\mid x) = \mathcal{N}(\mu, \operatorname{diag}\sigma^2)$:
\begin{equation}
  \mathcal{L}_\mathrm{KL} = \tfrac12 \textstyle\sum \bigl(\mu^2 + \sigma^2 - \log\sigma^2 - 1\bigr),
\end{equation}
summed over all latent entries and averaged over the batch. It is added with
weight $w_\mathrm{KL}\in\{10^{-7},10^{-5}\}$.

\paragraph{Score distillation (SDS).} A frozen DiT $v_\phi$ acts as a critic on
the latents $z = E_\theta(x)$ of the VAE being trained. With
$\sigma$ drawn from the shifted logit-normal distribution used in DiT training
and $\varepsilon\sim\mathcal{N}(0,I)$,
\begin{equation}
  z_\sigma = (1-\sigma)\,z + \sigma\,\varepsilon, \qquad
  \mathcal{L}_\mathrm{SDS} = \bigl\lVert v_\phi(z_\sigma;\sigma,\gamma) - (\varepsilon - z)\bigr\rVert^2 .
\end{equation}
The first latent frame is kept clean, as in DiT training, and is excluded from
the loss. The critic's forward pass runs without gradients, so the gradient
reaches the encoder only through the target $\varepsilon - z$. The critic is
the stock pretrained LTX-Video 2B transformer, never fine-tuned on this data,
with a randomly initialised $\gamma$ embedding. The term is added with weight
$w_\mathrm{SDS}\in\{10^{-4},10^{-2},1\}$.

\begin{table}[htbp]
\centering
\small
\caption{The single setting each arm changes relative to the baseline.}
\label{tab:app_arms}
\begin{tabular}{lll}
\toprule
Group & Arm & Change relative to the baseline \\
\midrule
Adaptation & Unfinetuned & no training (widened VAE with zero-initialised pressure slot) \\
           & Colour adapter & VAE encoder frozen; trainable decoder and $1\times1$ in/out adapters (4$\leftrightarrow$3 ch., 32 hidden, GELU, identity-initialised) \\
           & Decoder only & encoder frozen, including \texttt{conv\_in} \\
           & Decoder + \texttt{conv\_in} & encoder frozen except \texttt{conv\_in} \\
\midrule
Loss & RMSE only & $w_{H^1} = 0$ \\
     & RMSE + H1 + SSIM & $w_\mathrm{SSIM} = 0.15$ added \\
\midrule
Learning rate & LR $10^{-5}$ & peak learning rate $10^{-5}$ \\
\midrule
Resolution & Trained at $128^2$ & $128\times128$ data (and its normalisation statistics) \\
\midrule
Latent reg. & KL & $w_\mathrm{KL} \in \{10^{-7}, 10^{-5}\}$ \\
            & SDS & $w_\mathrm{SDS} \in \{10^{-4}, 10^{-2}, 1\}$ \\
\bottomrule
\end{tabular}
\end{table}

Two further arms were started and excluded because training collapsed in the
second epoch, right after warm-up reached the peak learning rate: a learning
rate of $10^{-4}$ and $w_\mathrm{KL}=10^{-6}$. The collapse of $w_\mathrm{KL}=10^{-6}$
is treated as a training instability rather than an effect of the KL term,
because the two KL weights on either side of it trained normally.

\section{DiT fine-tuning and inference}
\label{app:setup_dit}

For each VAE, the training and validation simulations are encoded once with the
posterior mean, and the full transformer is fine-tuned on these latents with a
flow-matching objective (no low-rank adapters).

\begin{table}[htbp]
\centering
\small
\caption{DiT fine-tuning and inference configuration, identical for every VAE.}
\label{tab:app_dit_hparams}
\begin{tabular}{ll}
\toprule
Setting & Value \\
\midrule
Trainable modules & full LTX-Video 2B transformer and the $\gamma$ embedding \\
Objective & flow matching, timesteps from a shifted logit-normal distribution \\
Conditioning & first latent frame kept clean ($p=1$); $\gamma$ as 4 tokens of width 4096 \\
$\gamma$ embedding & 64 Fourier features $\to$ MLP (hidden 256), dropout 0.1, $\gamma\in[1,2]$ \\
Optimiser & AdamW, learning rate $10^{-5}$, cosine decay to $10^{-8}$ \\
Steps & 8000 (batch size 32: 1 per device $\times$ 4 accumulation $\times$ 8 devices) \\
Gradient clipping & global norm 1.0 \\
Precision & bf16 mixed precision, gradient checkpointing \\
\midrule
Inference & 20 denoising steps, no classifier-free guidance (scale 1.0) \\
Rollout & frame 0 encoded with the posterior mean; 105 frames generated, trimmed to 101 \\
Seed & 42 (+ simulation index) \\
\bottomrule
\end{tabular}
\end{table}

\section{Hardware and run times}
\label{app:setup_hw}

Most runs used one node with four Intel Data Center GPU Max 1550 cards, exposed
as eight devices (one per tile). The Decoder + \texttt{conv\_in} arm was trained
on one node with four NVIDIA GH200 GPUs, at the same effective batch size of 32
(4 per GPU, 2 accumulation steps). Table~\ref{tab:app_runtime} lists typical
wall-clock times.

\begin{table}[htbp]
\centering
\small
\caption{Typical wall-clock times.}
\label{tab:app_runtime}
\begin{tabular}{llr}
\toprule
Job & Hardware & Time \\
\midrule
VAE, full fine-tuning, $256^2$, 20 epochs & $4\times$ Max 1550 & 15.7--16.7\,h \\
VAE, decoder only, $256^2$ & $4\times$ Max 1550 & 8.6\,h \\
VAE, full fine-tuning, $128^2$ & $4\times$ Max 1550 & 7.1\,h \\
VAE, full fine-tuning, $256^2$ & $4\times$ GH200 & $\approx$3.5\,h \\
DiT, 8000 steps & $4\times$ Max 1550 & 19.5--20.1\,h \\
VAE + DiT evaluation, 500 test simulations & $4\times$ Max 1550 & 1.5\,h \\
\bottomrule
\end{tabular}
\end{table}

% ---------------------------------------------------------------------------
\chapter{Dataset and Evaluation Protocol}
\label{app:data}

\section{Data and splits}
\label{app:data_splits}

The experiments use the Euler multi-quadrant data set (\texttt{euler\_mq}):
two-dimensional compressible Euler simulations started from piecewise-constant
initial states and parameterised by the heat-capacity ratio $\gamma$. Each
simulation has 101 frames of four fields: density $\rho$, momentum $m_x$ and
$m_y$, and pressure $p$. The data set ships at $512^2$ with downsampled
$256^2$ and $128^2$ copies. All experiments use $256^2$, except the resolution
arm.

\begin{table}[htbp]
\centering
\small
\caption{Data splits. The validation split is a fixed, seeded permutation
(seed 42) of the training file; the test file is used only for final
evaluation.}
\label{tab:app_splits}
\begin{tabular}{lrl}
\toprule
Split & Simulations & Use \\
\midrule
Training & 4000 & VAE and DiT training \\
Validation & 500 & checkpoint selection, training curves \\
Test (separate file) & 500 & all reported results \\
\bottomrule
\end{tabular}
\end{table}

The test set contains ten values of $\gamma$ (1.13, 1.22, 1.30, 1.33, 1.365,
1.40, 1.404, 1.453, 1.597, 1.76), with 50 simulations each. The training file
contains the same values, with 450 simulations each.

\section{Normalisation}
\label{app:data_norm}

Each channel $c$ is standardised with the mean $\mu_c$ and standard deviation
$s_c$ of the training data, then scaled and clipped to $[-1,1]$:
\begin{equation}
  \tilde u_c = \operatorname{clip}\!\left(\frac{u_c - \mu_c}{5\,s_c},\,-1,\,1\right).
\end{equation}
At $256^2$, $\mu = (0.905, 0.004, -0.002, 0.903)$ and
$s = (0.662, 0.534, 0.527, 0.805)$ for $(\rho, m_x, m_y, p)$. The $128^2$
statistics differ by less than $0.2\%$. All losses and metrics are computed on
the normalised fields; figures show physical units.

\section{Metrics}
\label{app:data_metrics}

For a prediction $\hat u$ and target $u$ of one simulation, the
variance-normalised RMSE is
\begin{equation}
  \mathrm{VRMSE}(\hat u, u) = \sqrt{\frac{\operatorname{mean}\bigl((\hat u - u)^2\bigr)}{\operatorname{var}(u)}} ,
\end{equation}
with mean and variance over all frames and pixels. Three variants are reported:
\begin{itemize}
  \item \emph{per-channel VRMSE}: computed separately for each field;
  \item \emph{channel-mean VRMSE}: the mean of the four per-channel values. This
        is the headline metric, because it weights every field equally
        regardless of its variance;
  \item \emph{pooled VRMSE}: all four channels treated as one tensor.
\end{itemize}
Each metric is computed per simulation and then averaged over the 500 test
simulations. The standard error of the mean (SEM) is taken over these 500
values. All arms are evaluated on the same simulations, so two arms can also be
compared per simulation.

The DiT evaluation runs two passes per simulation. \emph{VAE only} encodes and
decodes the full ground-truth sequence, which gives the reconstruction floor.
\emph{VAE + DiT} encodes only frame~0, lets the DiT generate the remaining
frames from it and $\gamma$, and decodes the result. The latent VRMSE compares
the generated latent with the encoder's latent of the ground truth, and so
measures the DiT's error without the decoder.

\section{Run-to-run variation}
\label{app:data_noise}

Repeating the baseline with three seeds at $128^2$ gave a spread of about
$1.1\%$ in validation VRMSE. At $256^2$, the baseline was trained twice with the
same seed on two different clusters (eight Max 1550 tiles vs.\ four GH200 GPUs,
which changes data sharding and numerics). The two runs differ by $0.8\%$ in
channel-mean validation VRMSE. This is a single pair, not a seed study.
Differences of about $1\%$ or less between arms are therefore treated as ties.
The standard errors of Appendix~\ref{app:results} capture only the sampling of
test simulations, not this training variation.

% ---------------------------------------------------------------------------
\chapter{Complete Results}
\label{app:results}

Table~\ref{tab:app_vae_test} gives the full VAE reconstruction results on the
test set, and Table~\ref{tab:app_vae_val} the same arms on the validation split.
Validation and test agree closely. The ranking of the arms is the same on both
sets, and the relative change of each arm differs by at most 0.6 percentage
points between them. The exceptions are the two arms far from the baseline:
Decoder only ($+77.4\%$ vs.\ $+80.0\%$) and the colour adapter ($+493\%$ vs.\ $+480\%$).
Table~\ref{tab:app_dit} gives the DiT results per VAE.

% generated by scripts/thesis_tables.py -- do not edit by hand
\begin{table}[htbp]
\centering
\small
\setlength{\tabcolsep}{4pt}
\caption{VAE reconstruction VRMSE on the held-out test set (500 simulations, $256\times256$). $\Delta$ is the change of the channel-mean VRMSE relative to the baseline. Lower is better.}
\label{tab:app_vae_test}
\begin{tabular}{lrrrrrrr}
\toprule
Arm & Pooled & Ch.\ mean $\pm$ SEM & $\Delta$ & $\rho$ & $m_x$ & $m_y$ & $p$ \\
\midrule
\multicolumn{8}{l}{\textit{Adaptation}} \\
Unfinetuned (inflated only) & 0.4774 & 0.5569 $\pm$ 0.0117 & +595.9\% & 0.1269 & 0.1116 & 0.1302 & 1.8589 \\
Colour adapter & 0.4081 & 0.4643 $\pm$ 0.0102 & +480.2\% & 0.1050 & 0.0863 & 0.0924 & 1.5736 \\
Decoder only & 0.1061 & 0.1441 $\pm$ 0.0023 & +80.0\% & 0.0974 & 0.0870 & 0.0938 & 0.2980 \\
Decoder + conv\_in & 0.0677 & 0.0929 $\pm$ 0.0015 & +16.1\% & 0.0843 & 0.0772 & 0.0803 & 0.1299 \\
\textbf{Full FT (baseline)} & 0.0590 & \textbf{0.0800 $\pm$ 0.0014} & -- & 0.0727 & 0.0707 & 0.0704 & 0.1064 \\
\midrule
\multicolumn{8}{l}{\textit{Loss}} \\
RMSE only & 0.0603 & 0.0814 $\pm$ 0.0014 & +1.7\% & 0.0767 & 0.0718 & 0.0724 & 0.1048 \\
RMSE + H1 + SSIM (0.15) & 0.0656 & 0.0889 $\pm$ 0.0015 & +11.1\% & 0.0808 & 0.0787 & 0.0786 & 0.1174 \\
\midrule
\multicolumn{8}{l}{\textit{Learning rate}} \\
LR $1\times10^{-5}$ & 0.0683 & 0.0938 $\pm$ 0.0015 & +17.2\% & 0.0805 & 0.0773 & 0.0778 & 0.1395 \\
\midrule
\multicolumn{8}{l}{\textit{Resolution}} \\
Trained at $128^2$ & 0.0882 & 0.1200 $\pm$ 0.0019 & +49.9\% & 0.1202 & 0.1009 & 0.1016 & 0.1572 \\
\midrule
\multicolumn{8}{l}{\textit{Latent reg.}} \\
KL $10^{-7}$ & 0.0593 & 0.0804 $\pm$ 0.0014 & +0.5\% & 0.0728 & 0.0707 & 0.0705 & 0.1077 \\
KL $10^{-5}$ & 0.0650 & 0.0877 $\pm$ 0.0014 & +9.6\% & 0.0807 & 0.0792 & 0.0784 & 0.1126 \\
SDS $10^{-4}$ & 0.0585 & 0.0793 $\pm$ 0.0014 & $-$0.9\% & 0.0723 & 0.0703 & 0.0700 & 0.1046 \\
SDS $10^{-2}$ & 0.0595 & 0.0806 $\pm$ 0.0014 & +0.8\% & 0.0734 & 0.0714 & 0.0709 & 0.1069 \\
SDS $1$ & 0.0766 & 0.1040 $\pm$ 0.0015 & +29.9\% & 0.0963 & 0.0941 & 0.0940 & 0.1316 \\
\bottomrule
\end{tabular}
\end{table}

% generated by scripts/thesis_tables.py -- do not edit by hand
\begin{table}[htbp]
\centering
\small
\setlength{\tabcolsep}{4pt}
\caption{VAE reconstruction VRMSE on the validation split (500 simulations held out from the training file) at the checkpoint with the best validation VRMSE. $\Delta_\mathrm{val}$ and $\Delta_\mathrm{test}$ compare each arm with the baseline on the validation and test sets. The model trained at $128^2$ is validated on $128^2$ data.}
\label{tab:app_vae_val}
\begin{tabular}{lrrrrrrrrr}
\toprule
Arm & Ep. & Pooled & Ch.\ mean & $\Delta_\mathrm{val}$ & $\Delta_\mathrm{test}$ & $\rho$ & $m_x$ & $m_y$ & $p$ \\
\midrule
\multicolumn{10}{l}{\textit{Adaptation}} \\
Colour adapter & 20 & 0.4042 & 0.4691 & +493.0\% & +480.2\% & 0.1031 & 0.0846 & 0.0921 & 1.5967 \\
Decoder only & 20 & 0.1013 & 0.1403 & +77.4\% & +80.0\% & 0.0956 & 0.0867 & 0.0938 & 0.2853 \\
Decoder + conv\_in & 20 & 0.0658 & 0.0918 & +16.0\% & +16.1\% & 0.0821 & 0.0766 & 0.0804 & 0.1280 \\
Full FT (baseline) & 20 & 0.0575 & 0.0791 & -- & -- & 0.0709 & 0.0700 & 0.0705 & 0.1051 \\
\midrule
\multicolumn{10}{l}{\textit{Loss}} \\
RMSE only & 20 & 0.0587 & 0.0805 & +1.8\% & +1.7\% & 0.0748 & 0.0713 & 0.0726 & 0.1036 \\
RMSE + H1 + SSIM (0.15) & 20 & 0.0639 & 0.0879 & +11.1\% & +11.1\% & 0.0787 & 0.0780 & 0.0787 & 0.1161 \\
\midrule
\multicolumn{10}{l}{\textit{Learning rate}} \\
LR $1\times10^{-5}$ & 20 & 0.0666 & 0.0928 & +17.3\% & +17.2\% & 0.0785 & 0.0765 & 0.0780 & 0.1382 \\
\midrule
\multicolumn{10}{l}{\textit{Resolution}} \\
Trained at $128^2$ & 20 & 0.0871 & 0.1181 & +49.3\% & +49.9\% & 0.1096 & 0.1083 & 0.1085 & 0.1461 \\
\midrule
\multicolumn{10}{l}{\textit{Latent reg.}} \\
KL $10^{-7}$ & 20 & 0.0578 & 0.0795 & +0.5\% & +0.5\% & 0.0710 & 0.0700 & 0.0707 & 0.1064 \\
KL $10^{-5}$ & 20 & 0.0633 & 0.0868 & +9.8\% & +9.6\% & 0.0788 & 0.0787 & 0.0785 & 0.1113 \\
SDS $10^{-4}$ & 20 & 0.0570 & 0.0784 & $-$0.9\% & $-$0.9\% & 0.0705 & 0.0696 & 0.0702 & 0.1034 \\
SDS $10^{-2}$ & 20 & 0.0580 & 0.0797 & +0.8\% & +0.8\% & 0.0717 & 0.0708 & 0.0710 & 0.1056 \\
SDS $1$ & 20 & 0.0746 & 0.1029 & +30.0\% & +29.9\% & 0.0943 & 0.0933 & 0.0939 & 0.1299 \\
\bottomrule
\end{tabular}
\end{table}

% generated by scripts/thesis_tables.py -- do not edit by hand
\begin{table}[htbp]
\centering
\small
\setlength{\tabcolsep}{4pt}
\caption{Channel-mean VRMSE of the VAE reconstruction alone and of the full pipeline (encode frame 0, DiT rollout, decode) on the test set (500 simulations), for a DiT trained for 8k steps on each VAE's latents. $\Delta$ is relative to the baseline VAE; the per-channel columns are the VAE + DiT errors; \emph{Latent} is the VRMSE of the rolled-out latent against the encoder's latent of the ground truth.}
\label{tab:app_dit}
\begin{tabular}{lrrrrrrrr}
\toprule
VAE & VAE only & VAE + DiT $\pm$ SEM & $\Delta$ & $\rho$ & $m_x$ & $m_y$ & $p$ & Latent \\
\midrule
\textbf{Full FT (baseline)} & 0.0802 & \textbf{0.4984 $\pm$ 0.0060} & -- & 0.514 & 0.506 & 0.516 & 0.459 & 0.558 \\
Unfinetuned (inflated only) & 0.5579 & 0.9501 $\pm$ 0.0128 & +90.6\% & 0.637 & 0.646 & 0.657 & 1.859 & 0.583 \\
KL $10^{-7}$ & 0.0807 & 0.4891 $\pm$ 0.0059 & $-$1.9\% & 0.507 & 0.499 & 0.506 & 0.445 & 0.558 \\
KL $10^{-5}$ & 0.0879 & 0.4733 $\pm$ 0.0053 & $-$5.0\% & 0.493 & 0.480 & 0.488 & 0.433 & 0.273 \\
SDS $10^{-4}$ & 0.0795 & 0.4974 $\pm$ 0.0060 & $-$0.2\% & 0.512 & 0.504 & 0.514 & 0.459 & 0.556 \\
SDS $10^{-2}$ & 0.0808 & 0.4583 $\pm$ 0.0054 & $-$8.1\% & 0.483 & 0.465 & 0.474 & 0.411 & 0.290 \\
SDS $1$ & 0.1044 & 0.5678 $\pm$ 0.0051 & +13.9\% & 0.572 & 0.578 & 0.585 & 0.536 & 0.065 \\
\bottomrule
\end{tabular}
\end{table}

